%% Introduction, input by paper.tex. Follows the order of abstract-v2.tex
%% and notes/intro-outline.md. Numbers come from Table 1 and Table 2
%% (loc.py --mu) and will change; update them before submission. Prior
%% work is cited by author name (see notes/related-work.md for the
%% evidence behind each statement). Siek's earlier work is cited in the
%% third person for double-blind review.

\section{Introduction}
\label{sec:intro}

This paper shows that denotational semantics can carry a proof of
compiler correctness through closure conversion, and it compares that
approach with operational ones, pass by pass. We give the first denotational proofs that closure conversion
both preserves and reflects the meaning of programs, divergence
included, for a Turing-complete language, one proof in a graph model
and one in Scott domains. We extend both proofs to an eight-pass
compiler and mechanize them in Agda. We then compare them with three
mechanized operational proofs of the same compiler, one each with
clocked interpreters, abstract machines, and reduction semantics. No style gives the smallest
proof for every pass. Overall, the proof with clocked interpreters is
the smallest and the graph-model proof is the second smallest. The
denotational proofs are short when the correctness of a pass is an
equation between meanings and long when a pass needs a logical
relation, whereas the operational proofs tend to grow with how much
the runs of the source and target differ. Section~\ref{sec:compiler} presents
the compiler, Section~\ref{sec:graph-models} the graph model, and
Section~\ref{sec:compiler-proof} the proofs.
Section~\ref{sec:evaluation} compares the five proofs,
Section~\ref{sec:related} discusses related work, and
Section~\ref{sec:conclusion} concludes. The appendices give the
semantics of every language, define every pass, and give the
Scott-domain proof of \pass{split closure}.

Anyone who sets out to prove a compiler correct must first choose a
style of semantics. CompCert \citep{Leroy:2009aa} gives its languages
small-step semantics and proves each pass by a simulation. CakeML
\citep{Kumar:2014aa} gives its languages clocked interpreters, which
\citet{Owens:2016aa} call functional big-step semantics, and proves
each pass by induction on evaluation. Denotational semantics offers a
different recipe. The meaning of a program is a mathematical object, a
compiler pass is correct when the meaning of a program before the pass
equals its meaning afterward, and the proof is by equational
reasoning. Early compiler proofs followed this recipe
\citep{Polak:1981aa,Clinger:1984aa}. We ask whether it extends to a
compiler for a higher-order language, and how its proofs compare in size.

However, the recipe meets an obstacle. Every compiler for a language
with lexically scoped functions performs \emph{closure conversion},
which replaces each function with a pair of code and the values of its
free variables, and replaces each call with code that takes the pair
apart. Closure conversion turns functions into data, so its
correctness proof must relate two representations of a function. No
proof of compiler correctness with respect to denotational semantics
covers closure conversion for a Turing-complete source language, and
the closest proofs each stop short. \citet{Chlipala:2007fk} proved
correct a compiler from the simply typed $\lambda$-calculus to
assembly, including closure conversion, but every program of that
language terminates, so there is no divergence to preserve.
\citet{Nielsen:2000aa} used a domain-theoretic semantics to prove that
defunctionalization, a close cousin of closure conversion, preserves
the meaning of programs with recursive functions, and left the other
direction for later, writing ``Eventually, we want to prove full correctness,
i.e., semantic equivalence of the two programs.''
\citet{Banerjee:2001aa} proved both directions, but gave the source
language an operational semantics. They report that they ``know of no
other proofs of correctness (that is, both directions) for ordinary or
CFA-based defunctionalization in the presence of recursion.'' Other
denotational proofs avoid closure conversion, because their source
language is first-order
\citep{Polak:1981aa,Xia:2019aa,Cheng:2026aa} or their target language
has closures \citep{Clinger:1984aa,Benton:2009aa}. When
\citet{Chlipala:2010aa} proved correct a compiler for a language with
recursion, he switched to operational semantics, explaining that the
denotational approach of his 2007 compiler ``cannot be applied
directly to Turing-complete programming languages.''

A correct compiler must both preserve and reflect the behavior of
programs. \emph{Semantic preservation} says that if the source program
produces an output, the target program produces the same output, and
\emph{semantic reflection} says that if the target program produces an
output, the source program produces the same output. By
contraposition, reflection guarantees that a diverging source program
compiles to a diverging target program. For closure conversion,
reflection is the hard direction for a denotational proof, because the
relation between source functions and target closures refers to itself
in negative position, and the domain of meanings is recursive. In a
graph model there is a further problem. The meaning of a function's
code describes its behavior on every environment, not just the
environment with which the compiled program pairs it. The target
denotation therefore contains code--environment combinations that
correspond to no source function, and a reflection proof must show
that these combinations never contribute an answer.

These observations lead to three research questions:
\begin{description}
\item[RQ1] Can a denotational semantics prove a compiler correct for a
  Turing-complete language with lexically scoped functions, through
  closure conversion, both preserving and reflecting the meaning of
  programs?
\item[RQ2] How do the sizes of mechanized correctness proofs of one
  compiler compare across five semantic styles, two denotational and
  three operational, pass by pass and overall?
\item[RQ3] Which properties of a pass determine the size of
  its correctness proof in each style?
\end{description}

We answer RQ1 affirmatively, in two models. Our first model is a graph
model, in which a program denotes a set of inductively defined
\emph{graph elements}
\citep{Engeler:1981aa,Plotkin:1993ab,Scott:1976lq}. Each element is a
finite piece of evidence about the program's behavior, the meaning of
a function is a set of input--output entries, and a program that
diverges produces no evidence, so its meaning is the empty set. For
reflection we introduce a new kind of relation between denotations,
the \emph{approximant logical relation}. Our key insight is to relate
a finite set of target graph elements (a finite approximation) to a
possibly infinite source denotation, and to define the relation by
recursion on the depth of the elements in the finite set
(Section~\ref{sec:approx-lr}). Our second model uses Scott domains, in
which a program denotes an element of a universal domain of tagged
values and divergence is the least element. There we define the
relation by recursion on the stages of the bilimit that constructs the
domain, as \citet{Reynolds:1974aa} did to relate direct and
continuation semantics. At each stage, the relation relates an
approximation of one value to the whole other value
(Section~\ref{sec:scott-delay}).

We use these results in two correctness proofs of an eight-pass
compiler. Its source language is \fpc, the call-by-value
$\lambda$-calculus with sums, products, and iso-recursive types of
\citet{Plotkin:1985aa}, extended with integers. Recursive types let us
write the call-by-value fixed-point combinator, so \fpc\ has general
recursion and admits divergence. The compiler first applies a
source-to-source optimizer modeled on the cp0 pass of Chez Scheme
\citep{Waddell:1997fk}, then performs closure conversion in three
passes, and finally uses four passes to remove complex operands, lift
locals, make control flow explicit, and select instructions of
\instr, a typed low-level intermediate language. Every language in
the compiler has a semantics in each model. For seven of the eight
passes, the proof in either model is an equation between the meanings
before and after the pass. The exception is the pass that separates
the code of each closure from its environment, \pass{split closure},
whose reflection proof uses the approximant logical relation in the
graph model and the stage relation in Scott domains. The
whole-compiler theorem says that the source and target denotations
contain the same constants and are nonempty together. We also prove
that the denotational semantics of \fpc\ and of \instr\ are adequate
with respect to their interpreters, which makes our theorem equivalent
to the usual operational one. We mechanize all of these results in
Agda.

To answer RQ2 and RQ3, we compare the two denotational proofs with
three operational proofs of the same compiler, each against its own
semantics: clocked interpreters with relations on values, abstract
machines with simulations, and reduction semantics with simulations.
Earlier comparisons of semantic styles for compiler proofs cover a
single phase or pass \citep{Owens:2016aa,Kahn:2025aa}; ours is the
first on a whole compiler for a higher-order language. We leave other
styles to future work, among them synthetic domain theory
\citep{Rosolini:1986aa,Hyland:1991aa} and its guarded variant
\citep{Birkedal:2011aa,Mogelberg:2016aa}, categorical semantics
\citep{Fiore:1994aa}, coinductive and pretty-big-step semantics
\citep{Leroy:2009rc,Chargueraud:2013aa}, the partiality monad
\citep{Danielsson:2012aa}, game semantics \citep{Hyland:2000ve}, and
action semantics \citep{Mosses:1996rp}. We did not use interaction
trees \citep{Xia:2019aa}, because a direct denotation of a
higher-order function as an interaction tree needs a type of values
that occurs negatively in its own definition, and the positivity
checks of Coq and Agda reject such a definition. Guarded interaction trees
\citep{Frumin:2024aa} lift this restriction by moving to guarded type
theory.

No style gives the smallest proof for every pass (RQ2). Counting each
proof against its own semantics, the clocked interpreters give the
smallest proof, at 1,504 lines of Agda, and the graph model gives the
second smallest, at 1,808, ahead of abstract machines, Scott domains, and
reduction semantics. If we also count the language-independent
infrastructure, the graph model ranks third and Scott domains fifth.
Two factors largely explain the sizes (RQ3). The size of a
denotational proof tends to depend most on whether the pass preserves the meaning of a program
exactly. The source optimizer, \pass{simplify}, does, and its
graph-model and Scott-domain proofs take 349 and 411 lines, against
508 to 697 for the operational proofs (Section~\ref{sec:equational}).
\pass{split closure} does not, and its proofs take 450 and 723 lines,
against 35 to 126. The size of an operational proof tends to depend
instead on how closely the run of the target follows the run of the source.
The runs correspond closely for \pass{split closure}, which is why its
operational proofs are small, and differ more for \pass{simplify},
which inlines functions. Without \pass{split closure}, the graph-model
proof would be the smallest of the five. The Scott-domain proof is 1.4
times as large as the graph-model proof. Much of the difference is in
\pass{split closure}, where we must show that the stage relation is
admissible and preserved by the maps between stages, and where the
method of \citet{Pitts:1996aa} does not
apply directly (Section~\ref{sec:evaluation}).

Our contributions are:
%% The technical contributions, at the end of the introduction, input by
%% paper.tex. One sentence each (two in a pinch); each is novel (to our knowledge, per
%% notes/novelty-literature.md: C2-C4, C6, C7, C9) and is expanded in the
%% section it points to. The section labels are placeholders until those
%% sections exist.

Our technical contributions are:
\begin{itemize}
\item A proof, mechanized in Agda, that a ten-pass compiler from a
  call-by-value $\lambda$-calculus with iso-recursive types, through
  closure conversion, to a typed assembly language is correct, using a
  graph-model semantics for every intermediate language. Most passes are
  proved by an equation between the denotations of a program before and
  after the pass (\S\ref{sec:compiler-proof}).
\item Approximant logical relations, which relate two denotational
  semantics by recursion on the depth of a finite approximation of one
  of them, and with which we prove the reflection direction of lowering
  closures into code and environment pairs, where recursive types rule
  out recursion on types (\S\ref{sec:approx-lr}).
\item Proofs that the graph-model semantics of the source language and of
  the assembly language are adequate with respect to interpreters for
  those languages, which yield an end-to-end theorem that the source
  program and the compiled program return the same integer or both
  diverge (\S\ref{sec:adequacy}).
\item A comparison with proofs of the same compiler in three operational
  styles (abstract machines, clocked interpreters and reduction
  semantics), which shows that the denotational proofs are among the
  smallest when program values stay at the same level of abstraction and
  are the largest when values are lowered (\S\ref{sec:evaluation}).
\end{itemize}

